% ============================================================================
%  Capitolo 0 -- Introduzione al corso                           (Lezione 00)
% ============================================================================
\chapter{Introduzione: prima programmare, poi formalizzare}
\label{ch:intro}

\epigraphline{Matrix è un sistema, Neo. E quel sistema è nostro nemico.}
  {\emph{Matrix}, Lana e Lilly Wachowski, 1999}

\begin{chapgoals}
\begin{itemize}
\item Che cos'è il corso, chi lo tiene e come sono usate le sue novantasei ore.
\item Il metodo: prima viene un programma che gira, poi la sua teoria.
\item La macchina: \prg, \riscv e il debugger \prgqt.
\item Il progetto: \vermone, un gioco scritto da un file vuoto e pubblicato.
\item Le regole: programma, risultati di apprendimento, valutazione, materiale.
\end{itemize}
\end{chapgoals}
% Generated by tools/textbook/export_listings.js from tools/decks/mastery.js and tools/concepts/registry.json: do not edit.
\begin{chapstatus}{Stato di apprendimento di questo capitolo}
\masterychip{now}{DA COMPRENDERE OGGI}\enspace il metodo: prima il programma; le regole del corso.\par\smallskip
\masterychip{use}{DA UTILIZZARE OGGI}\enspace il programma completo, gli strumenti di PRG32 (si usa, non si studia, in questo corso).\par\smallskip
\masterychip{later}{DA COMPRENDERE PIÙ AVANTI}\enspace i componenti di PRG32 (Capitolo~\ref{ch:setup}); RISC e un ISA aperto (Capitolo~\ref{ch:history}); il progetto Verm1 (Capitolo~\ref{ch:verm1p1}).
\end{chapstatus}
\begin{chapstatus}{Percorso del capitolo}
Alla fine del capitolo saprete:
\begin{enumerate}[nosep,leftmargin=2em]
\item descrivere i quattro passi della spirale di apprendimento e dire perché la previsione viene prima dell'esecuzione
\item enunciare le regole di valutazione: una prova intercorso, due parti d'esame da superare entrambe, nessun dispositivo e nessuna IA
\end{enumerate}
Si appoggia su: niente: è il primo capitolo. Prossimo passo: Capitolo~\ref{ch:setup}.
\end{chapstatus}


Un calcolatore è l'oggetto più complicato che la maggior parte delle persone
usa ogni giorno, e quello di cui meno sono incoraggiate a guardare il
funzionamento interno. I suoi costruttori hanno passato settant'anni a
nascondere la macchina dietro strati ciascuno più facile da usare di quello
sottostante: le porte logiche dietro le istruzioni, le istruzioni dietro i
linguaggi, i linguaggi dietro le librerie, le librerie dietro le applicazioni,
e ora le applicazioni dietro gli assistenti conversazionali. Sono gli strati a
rendere utile l'informatica. Sono anche un sistema nel senso del film citato
qui sopra: ci si può vivere dentro per una vita senza sapere che sotto c'è
qualcosa.

L'architettura dei calcolatori è la disciplina che guarda sotto. Studia come
una macchina capace soltanto di accendere e spegnere correnti arrivi a
rappresentare numeri, testo, immagini e suono; come una manciata di operazioni
logiche, ripetute miliardi di volte al secondo, diventino aritmetica, decisioni
e memoria; e come il tutto sia organizzato in modo che una sequenza di ordini
semplicissimi, un \term{programma}, possa farle fare qualunque cosa sia
calcolabile. Questo capitolo spiega come il corso affronta questo studio, con
quale macchina e attraverso quale progetto, e poi enuncia le regole pratiche.

% ----------------------------------------------------------------------------
\section{Il corso in breve}

\emph{Architettura dei Calcolatori e Laboratorio di Architettura dei
Calcolatori}, codice ARC12, è un insegnamento del primo anno, primo semestre,
del Corso di Laurea in Informatica dell'Università degli Studi di Napoli
«Parthenope». Vale 12~CFU, appartiene al settore scientifico-disciplinare
INFO-01/A (Informatica) e consiste di 96~ore, divise in parti uguali tra
lezioni e laboratorio (Tabella~\ref{tab:intro:glance}).

\begin{table}[htb]
\centering\small
\caption{Il corso in breve.}\label{tab:intro:glance}
\begin{tabularx}{\textwidth}{@{}lL@{}}
\toprule
Nome & Architettura dei Calcolatori e Laboratorio di Architettura dei Calcolatori (ARC12)\\
Collocazione & Primo anno, primo semestre, Corso di Laurea in Informatica\\
Crediti & 12 CFU, 96 ore; settore INFO-01/A (Informatica)\\
Lezioni & Raffaele Montella, 6 CFU, 48 ore\\
Laboratorio & Mariacarla Staffa, 6 CFU, 48 ore\\
Ricevimento & Stanza 420, martedì 14:00--16:00, su appuntamento; anche su
  Microsoft Teams (codice del team \code{c7rf33g})\\
Crediti aggiuntivi & Accenture SuperSkills! dà 2 CFU come ulteriori conoscenze
  (\url{https://richiediaccesso.uniparthenope.it})\\
Prerequisiti & Nessuno\\
\bottomrule
\end{tabularx}
\end{table}

Non ci sono prerequisiti formali. Non si presuppone alcuna esperienza
precedente di programmazione, di elettronica o di matematica oltre il programma
scolastico: l'unico vero prerequisito è la passione per l'informatica e per le
tecnologie che la rendono possibile.

Le due metà del corso sono due viste della stessa materia. Nelle
\textbf{lezioni} si presentano le idee: informazione, numeri, algebra di Boole,
reti logiche, memoria, processore, ingresso e uscita. Ogni argomento si apre
con un programma che si può eseguire e viene poi formalizzato con definizioni,
metodi e dimostrazioni; ogni lezione si chiude con domande a risposta multipla
nello stile della prova intercorso; e ogni quattro lezioni un tutorial aggiunge
la parte successiva di \vermone, il gioco del corso. Nel \textbf{laboratorio}
le stesse idee vengono esercitate: rappresentazione dei numeri e tecniche di
conversione; analisi, sintesi e minimizzazione di circuiti logici;
programmazione in assembly \riscv sulla piattaforma \prg; e, infine, il
progetto dello studente, una variante di \vermone o un gioco di sua invenzione.
Questo libro è il testo delle lezioni; i suoi esercizi sono il materiale
naturale del laboratorio.

% ----------------------------------------------------------------------------
\section{Un approccio nuovo}

\begin{prgconcept}[Prima programmare. Poi formalizzare.]
Fin dai primi giorni lo studente scrive assembly \riscv che gira su un
processore vero. Ogni argomento formale risponde poi a una domanda che il
programma dello studente ha già sollevato.
\end{prgconcept}

L'ordine tradizionale della materia è deduttivo: la rappresentazione, poi la
logica, poi i circuiti, poi il processore, poi, finalmente, un programma. È
l'ordine in cui la macchina viene \emph{costruita}, e ha l'eleganza di ogni
costruzione dai primi principi. Non è l'ordine in cui la macchina viene
\emph{capita}. Uno studente che incontra il complemento a due nella terza
settimana di un corso deduttivo non ha mai visto un numero negativo dentro un
calcolatore, e non ha alcuna ragione di interessarsi a come si scrive. Uno
studente che, la settimana prima, ha spostato un rettangolo verso sinistra
finché la sua coordinata, mostrata dal debugger, è passata da \hex{00000000} a
\hex{FFFFFFFF} ha una domanda, e il complemento a due è la risposta.

Il corso è quindi organizzato come una spirale, ripetuta per ogni argomento
(Figura~\ref{fig:intro:spiral}).

\begin{figure}[htb]
\centering
\begin{tikzpicture}[node distance=7mm]
  \node[blkacc,text width=27mm,minimum height=20mm] (a) {\textbf{Gioca}\\[2pt]\footnotesize si scrivono poche istruzioni, si compila una cartuccia, la si esegue};
  \node[blkamb,text width=27mm,minimum height=20mm,right=of a] (b) {\textbf{Osserva}\\[2pt]\footnotesize un numero fa il giro, un bit accende una lampada, un rettangolo si muove};
  \node[blkvio,text width=27mm,minimum height=20mm,right=of b] (c) {\textbf{Formalizza}\\[2pt]\footnotesize definizioni, teoremi e metodi spiegano perché};
  \node[blkgrn,text width=27mm,minimum height=20mm,right=of c] (d) {\textbf{Ritorna}\\[2pt]\footnotesize si prevede il risultato prima di eseguire};
  \draw[->,thick,prgink] (a) -- (b); \draw[->,thick,prgink] (b) -- (c); \draw[->,thick,prgink] (c) -- (d);
  \draw[->,thick,prgink,rounded corners=3mm] (d.south) -- ++(0,-7mm) -| node[pos=0.25,below,lbl]{argomento successivo} (a.south);
\end{tikzpicture}
\caption{La spirale di apprendimento, ripetuta per ogni argomento del corso.}
\label{fig:intro:spiral}
\end{figure}

Prima si \emph{gioca}: si scrivono poche istruzioni, si compila e si esegue una
cartuccia. Poi si \emph{osserva}: succede qualcosa che non ci si aspettava,
oppure che ci si aspettava e che ora si può vedere. Poi si \emph{formalizza}:
definizioni e teoremi spiegano esattamente perché. Infine si \emph{ritorna}
alla macchina, e la misura della comprensione è la capacità di prevedere un
risultato prima di eseguire il programma. Un voto è un'opinione; una racchetta
che esce dallo schermo è un fatto. Il corso usa i fatti, e poi li spiega.

\subsection{Gli obiettivi del corso}

Gli obiettivi sono quelli di ogni corso serio sulla materia; cambia solo il
percorso. Il corso

\begin{itemize}
\item presenta gli aspetti fondamentali dell'architettura dei calcolatori
  moderni;
\item usa, in laboratorio, i concetti di base necessari a programmarli in modo
  efficiente con linguaggi di basso livello: linguaggio macchina e assembly;
\item tratta con rigore gli strumenti con cui si costruisce un processore,
  cioè l'algebra di Boole con la sintesi di reti logiche combinatorie e
  sequenziali, e la rappresentazione binaria con l'aritmetica binaria;
\item descrive in dettaglio la struttura funzionale e fisica della CPU, della
  memoria e del sottosistema di ingresso/uscita;
\item sviluppa codice assembly che rende visibile come lavora davvero un
  processore e come interagiscono le unità fondamentali di un calcolatore.
\end{itemize}

% ----------------------------------------------------------------------------
\section{La macchina: PRG32 e RISC-V}

\subsection{Una console rétro programmata in assembly}

\prg è un runtime didattico open source per insegnare l'assembly \riscv e il C,
sviluppato all'Università degli Studi di Napoli «Parthenope»
\citep{prg32repo}. All'occhio è una piccola console per videogiochi del tipo
che si vendeva nei primi anni Ottanta (Figura~\ref{fig:intro:prg32}); per il
programmatore è un contratto volutamente sottile.

\begin{figure}[htb]
\centering
\begin{subfigure}[b]{0.48\textwidth}\centering
\includegraphics[width=\textwidth]{menu_main.png}
\caption{Il menu del firmware residente.}
\end{subfigure}\hfill
\begin{subfigure}[b]{0.48\textwidth}\centering
\includegraphics[width=\textwidth]{game_pong.png}
\caption{Pong: codice macchina \riscv nativo.}
\end{subfigure}
\caption{\prg, la macchina del corso.}
\label{fig:intro:prg32}
\end{figure}

\begin{itemize}
\item Il codice dello studente gira in modo \emph{nativo}, come codice macchina
  \riscv a 32 bit dell'insieme di istruzioni RV32IMAC. Non è interpretato da
  un simulatore didattico: le istruzioni scritte sono le istruzioni eseguite.
\item Il contratto consiste di tre funzioni scritte dallo studente:
  \code{init} gira una volta; \code{update} e \code{draw} girano una volta per
  frame.
\item I giochi viaggiano come piccole \term{cartucce}, file con estensione
  \fname{.prg32}.
\item La macchina offre un viewport di 320 per 200 pixel, un joystick con
  pochi pulsanti, e il suono: abbastanza piccola da tenerla tutta in mente,
  abbastanza ricca da mettere alla prova l'intero corso.
\end{itemize}

\subsection{La filosofia di PRG32}

Sei principi, presi dal libro di testo della piattaforma
\citep{prg32textbook}, danno forma sia a \prg sia al modo in cui questo libro
lo usa.

\begin{description}[style=nextline,leftmargin=1.5em]
\item[Abbastanza piccola da capirla.] Un processore, uno schermo di 320 per
  200 pixel, un joystick, un piccolo altoparlante. Nulla è nascosto: l'intera
  macchina si può tenere in mente.
\item[Vera, non simulata.] Le istruzioni scritte sono quelle che il chip
  esegue. Ciò che si impara qui è vero su ogni processore \riscv, da un
  sensore a un supercalcolatore.
\item[I limiti rendono creativi.] Pochi kilobyte e tre funzioni. I giochi
  degli anni Ottanta sono nati da limiti come questi: un vincolo è un invito.
\item[Una cosa alla volta.] Si comincia con funzioni che si limitano a
  ritornare. Si aggiunge una cosa, si esegue, si guarda. \code{update} cambia
  lo stato, \code{draw} lo mostra. Il suono viene per ultimo.
\item[Guardare dentro.] Quando qualcosa sorprende, si ferma il programma e si
  leggono registri e memoria. Il debugger fa parte del metodo, non è uno
  strumento d'emergenza.
\item[Fare, poi condividere.] Una cartuccia è un file che si può dare a un
  amico. Il Cartridge Store è lo scaffale su cui il lavoro finito incontra i
  suoi giocatori.
\end{description}

La piattaforma ha quindi due usi nel corso. È una \emph{palestra}: ogni
capitolo si apre con un breve programma che gira e si chiude tornando alla
macchina; un esercizio consiste di una previsione e di una verifica; e la
macchina non lusinga mai, perché una maschera sbagliata accende la lampada
sbagliata. Ed è un'\emph{officina}: una cartuccia vuota e uno schermo, su cui
regole, colori, suoni e livelli sono decisioni di progetto che spettano
all'autore. Il progetto del corso comincia guidato e finisce libero, e ciò che
è finito si può pubblicare. Gli studenti del primo anno di rado si aspettano
che un corso di architettura chieda loro un'idea; questo lo fa.

\subsection{PRG32-QT: il debugger del corso}

\prgqt è \prg su un portatile: un'applicazione che esegue le stesse cartucce
senza scheda e senza compilatore \citep{prg32qt}. In modalità debug apre la
macchina (Figura~\ref{fig:intro:debugger}): il programma può essere messo in
pausa, ripreso e fatto avanzare di un'istruzione alla volta; si possono
mettere breakpoint su qualunque istruzione di \code{init}, \code{update} o
\code{draw}; vengono mostrati i trentadue registri e 128~byte di memoria a
qualunque indirizzo; la velocità si può regolare tra 0,01 e 4 volte quella
reale; e dei contatori riportano i frame, le istruzioni eseguite e i frame
arrivati in ritardo.

\begin{figure}[htb]
\centering
\includegraphics[width=0.92\textwidth]{prg32qt_debugger.png}
\caption{\prgqt in modalità debug: a sinistra il gioco, a destra la macchina.}
\label{fig:intro:debugger}
\end{figure}

Il debugger si usa in ogni capitolo: per verificare una previsione, per trovare
un errore, per misurare. Fornisce anche la definizione operativa di
comprensione adottata in questo libro.

\begin{prgconcept}[Quando si è capito un programma?]
Un programma è capito quando si sa dire che cosa conterrà ciascun registro
\emph{prima} di premere \key{Step}.
\end{prgconcept}

\subsection{Perché RISC-V?}

La scelta del processore non è neutra, e quattro ragioni hanno portato a
\riscv.

\begin{description}[style=nextline,leftmargin=1.5em]
\item[È aperto.] \riscv è un insieme di istruzioni a standard aperto: chiunque
  può progettare, costruire e vendere un processore \riscv senza pagare una
  licenza per l'architettura \citep{riscvisa}.
\item[È piccolo.] L'insieme di istruzioni intero di base ha una quarantina di
  istruzioni e trentadue registri. Lo si può imparare tutto in un semestre, e
  l'Appendice~\ref{app:riscv} lo elenca tutto.
\item[È reale.] \ins{addi}, \ins{lw}, \ins{sw} e \ins{beqz} non sono un
  dialetto didattico. L'industria consegna core \riscv a miliardi, ed essi
  eseguono esattamente queste istruzioni.
\item[È europeo.] L'Europa considera \riscv una leva per la sovranità
  tecnologica: lo European Chips Act e progetti come DARE (Digital Autonomy
  with RISC-V in Europe, un progetto della EuroHPC Joint Undertaking)
  costruiscono processori su di esso \citep{eurohpcdare}.
\end{description}

\subsection{La mappa: ogni argomento si incontra due volte}

La Tabella~\ref{tab:intro:map} è il corso in una pagina. Ogni argomento formale
della colonna di sinistra si incontra prima sulla macchina, nella forma data
dalla colonna di destra, e solo dopo nella sua forma generale. Vale la pena di
tornare a questa tabella all'inizio di ogni parte del libro.

\begin{table}[htb]
\centering\small
\caption{Ogni argomento formale si incontra prima su \prg.}\label{tab:intro:map}
\begin{tabularx}{\textwidth}{@{}P{0.36\textwidth}L@{}}
\toprule
\textbf{Argomento formale} & \textbf{Incontrato prima su \prg come\ldots}\\
\midrule
Notazione posizionale, esadecimale & il colore \hex{07E0} e il valore mostrato dal debugger\\
Complemento a due, overflow & una posizione che va a sinistra oltre lo zero e mostra \hex{FFFFFFFF}\\
Virgola fissa e virgola mobile & una palla che si muove di 0,375 pixel per frame su un chip senza unità in virgola mobile\\
Algebra di Boole, maschere di bit & \code{andi t3, t0, 1} per verificare se \key{LEFT} è premuto\\
Elettronica, porte logiche & il pulsante che collega un piedino del chip a massa\\
Reti combinatorie, sommatori & che cosa fa davvero \ins{add} a trentadue bit\\
Reti sequenziali, registri & il contatore che sopravvive da un frame al successivo\\
CPU, datapath, ciclo dell'istruzione & l'esecuzione passo-passo di un'istruzione in \prgqt\\
Memoria e ingresso/uscita & \ins{lw} e \ins{sw}, lo stack, il joystick e il display\\
\bottomrule
\end{tabularx}
\end{table}

% ----------------------------------------------------------------------------
\section{Il progetto: Verm1}

\vermone, che si legge \emph{vermone}, è un omaggio a Snake: il gioco del
cabinato \emph{Blockade} (1976) e dei telefoni Nokia (1997). Un verme si muove
su una griglia, mangia, cresce, e non deve mordersi. Lo studente ne scrive ogni
riga in assembly \riscv, senza motore e senza librerie oltre alle chiamate di
\prg. Il gioco è abbastanza piccolo per uno studente del primo anno e
abbastanza ricco da richiedere ogni argomento del corso; alla fine è pronto per
il Cartridge Store con sopra il nome del suo autore. L'implementazione di
riferimento, con cui ogni lettore può confrontare il proprio programma, è un
contributo di Giuseppe Coviello (Appendice~\ref{app:verm1}).

Il progetto è il filo del libro. Si costruisce in quattro parti e una
pubblicazione (Tabella~\ref{tab:intro:verm1}, Figura~\ref{fig:intro:verm1}), e
ogni parte usa esattamente ciò che i quattro capitoli precedenti hanno
insegnato. Ogni parte finisce con un gioco che gira: il lettore non torna mai
a casa con un programma rotto.

\begin{table}[htb]
\centering\small
\caption{\vermone: quattro parti e una pubblicazione.}\label{tab:intro:verm1}
\begin{tabularx}{\textwidth}{@{}lP{0.20\textwidth}LP{0.22\textwidth}@{}}
\toprule
\textbf{Capitolo} & \textbf{Dopo la parte su} & \textbf{\vermone acquista} & \textbf{Con}\\
\midrule
\ref{ch:verm1p1}\quad Parte 1 & informazione e numeri & una griglia, una testa che si muove, un mondo senza bordi & colori, scorrimenti, virgola fissa\\
\ref{ch:verm1p2}\quad Parte 2 & logica e circuiti & la sterzata, il cibo, un punteggio & maschere, XOR, tabelle\\
\ref{ch:verm1p3}\quad Parte 3 & stato e memoria & un corpo che cresce, le collisioni, tre schermate & una macchina a stati, vettori\\
\ref{ch:verm1p4}\quad Parte 4 & la macchina & cifre, livelli, suono, disegno veloce & chiamate, divisione, il frame buffer\\
\ref{ch:store}\quad Lo Store & l'intero corso & un nome, un'icona, i crediti, un bundle & gli strumenti di \prg\\
\bottomrule
\end{tabularx}
\end{table}

\begin{figure}[htb]
\centering
\begin{subfigure}[b]{0.24\textwidth}\includegraphics[width=\textwidth]{verm1_p1.png}\caption{Parte 1: una testa}\end{subfigure}\hfill
\begin{subfigure}[b]{0.24\textwidth}\includegraphics[width=\textwidth]{verm1_p2.png}\caption{Parte 2: cibo, punteggio}\end{subfigure}\hfill
\begin{subfigure}[b]{0.24\textwidth}\includegraphics[width=\textwidth]{verm1_p3_play.png}\caption{Parte 3: un corpo}\end{subfigure}\hfill
\begin{subfigure}[b]{0.24\textwidth}\includegraphics[width=\textwidth]{verm1_p4_play.png}\caption{Parte 4: finito}\end{subfigure}
\caption{\vermone cresce con il lettore: lo stesso file, quattro sedute. Le
schermate sono prese dai programmi di riferimento dai test automatici del
materiale del corso.}
\label{fig:intro:verm1}
\end{figure}

Un capitolo tutorial obbedisce a quattro regole. È \emph{autonomo}: mostra
l'intero programma, e ogni riga è nuova, modificata, oppure ristampata
invariata. Procede \emph{passo per passo}: poche righe alla volta, e dopo ogni
passo il programma viene compilato, eseguito e confrontato con un risultato
dichiarato. Usa \emph{il debugger}: ogni tutorial ha una sessione in \prgqt,
con un breakpoint, un registro da osservare e una regione di memoria da
leggere. E il suo risultato appartiene al lettore: ogni tutorial finisce con
modifiche da fare in autonomia, perché il gioco che viene pubblicato non
dovrebbe essere identico a quello stampato qui.

% ----------------------------------------------------------------------------
\section{Il corso in blocchi}

L'ordine degli argomenti è stato scelto per l'apprendimento. Prima viene che
cosa significa un bit; poi come i bit si combinano; poi come si ricordano; poi
la macchina che fa tutte e tre le cose. Il processore è presentato prima dei
dispositivi che pilota. Dopo ogni blocco il progetto lo usa. L'ultima lezione è
uno sguardo in avanti, tenuta prima di una visita ai sistemi di calcolo del
Dipartimento (Tabella~\ref{tab:intro:blocks}).

\begin{table}[htb]
\centering\small
\caption{I blocchi del corso. I numeri sono quelli delle lezioni e dei
capitoli di questo libro.}\label{tab:intro:blocks}
\begin{tabularx}{\textwidth}{@{}lLP{0.28\textwidth}@{}}
\toprule
\textbf{Blocco} & \textbf{Lezioni} & \textbf{Poi}\\
\midrule
Partenza & 0 Introduzione; 1 \prg e \prgqt: installazione; 2 Primo contatto con \riscv & \\
Informazione & 3 Storia; 4 Informazione digitale; 5 Interi; 6 Reali & 7 \vermone, parte 1\\
Logica & 8 e 9 Algebra di Boole; 10 Elettronica; 11 Macchine combinatorie & 12 \vermone, parte 2\\
Stato e memoria & 13 Macchine sequenziali; 14 Registri e bus; 15 Memoria; 16 Memoria permanente & 17 \vermone, parte 3\\
La macchina & 18 Processore; 19 Istruzioni e interruzioni; 20 Scheda video; 21 Ingresso e uscita & 22 \vermone, parte 4\\
Pubblicazione & & 23 \vermone sul Cartridge Store\\
Prospettive & 24 Oltre un solo core: \riscv superscalare, HPC, GPU e FPGA & la visita ai sistemi di calcolo\\
\bottomrule
\end{tabularx}
\end{table}

% ----------------------------------------------------------------------------
\section{Programma e risultati di apprendimento}

\subsection{Programma}

\begin{description}[style=nextline,leftmargin=1.5em]
\item[Organizzazione del calcolatore.] Hardware e software; software
  applicativo e di sistema; il modello di von Neumann; principi di
  funzionamento di un calcolatore.
\item[Algebra di Boole e logica.] Funzioni booleane, tabelle di verità; reti
  combinatorie e loro minimizzazione; sommatori; macchine a stati finiti; reti
  sequenziali; flip-flop, registri.
\item[Rappresentazione dei dati.] Codifica; codici ridondanti, rilevazione e
  correzione degli errori; decodificatori, multiplexer; sistemi posizionali;
  naturali, interi, reali; aritmetica.
\item[Struttura di un calcolatore.] La CPU: unità di controllo, ALU, registri;
  il ciclo dell'istruzione; il datapath; la memoria centrale; il sottosistema
  di I/O; architetture microprogrammate.
\item[Assembly RISC-V.] Insieme di istruzioni; registri, indirizzamento;
  aritmetica e logica; salti; sottoprogrammi, convenzione di chiamata, stack;
  memoria e I/O.
\item[Progetto.] \vermone: un gioco in assembly \riscv su \prg, costruito in
  quattro parti guidate, poi variato, documentato e confezionato dallo
  studente o da un piccolo gruppo.
\end{description}

Il programma dettagliato è pubblicato sul sistema informativo dell'ateneo e sul
portale di e-learning, \url{https://elearning.uniparthenope.it}.

\subsection{Risultati di apprendimento}

I risultati sono enunciati, come d'uso nell'istruzione superiore europea,
secondo i cinque descrittori di Dublino.

\begin{description}[style=nextline,leftmargin=1.5em]
\item[Conoscenza e capacità di comprensione.] Una solida conoscenza del sistema
  binario, della logica digitale e dei principi di funzionamento dei
  calcolatori. Non si presuppone alcuna formazione informatica precedente.
\item[Capacità di applicare conoscenza e comprensione.] La capacità di
  applicare la teoria a esercizi e a programmi in assembly \riscv, e di
  riconoscere sia le capacità sia i limiti dei calcolatori su problemi reali.
\item[Autonomia di giudizio.] La capacità di scrivere programmi assembly che
  rappresentano ed elaborano in binario problemi reali, scegliendo
  rappresentazioni dei dati e strutture di controllo adeguate.
\item[Abilità comunicative.] La capacità di partecipare attivamente a gruppi
  che progettano, realizzano e documentano applicazioni assembly, usando una
  terminologia tecnica corretta.
\item[Capacità di apprendimento.] La capacità di collegare rappresentazione,
  reti logiche e processori in un unico quadro, e di aggiornarsi in autonomia
  da fonti tecniche affidabili.
\end{description}

% ----------------------------------------------------------------------------
\section{Organizzazione e valutazione}

Questa sezione riporta le regole dell'anno accademico 2026/27. Le date non
ancora fissate al momento della compilazione del libro vengono comunicate in
aula e sul portale di e-learning, che in caso di differenze prevalgono su
questo testo.

\subsection{Orario}

Da settembre a dicembre 2026 il corso si tiene quattro volte alla settimana in
Aula~1: lezioni il lunedì (14:00--16:00) e il mercoledì (09:00--11:00),
laboratorio il martedì e il venerdì (09:00--11:00). Un calendario interno con
l'effettiva distribuzione di lezioni e sessioni di laboratorio è condiviso con
gli studenti.

\subsection{Una prova intercorso}

C'è una sola prova intercorso, svolta in presenza e su carta in circa
45~minuti: \textbf{dieci domande a risposta multipla più due esercizi}. Si
tiene di lunedì, 14:00--16:00, in Aula~1, in una data che sarà comunicata; gli
studenti possono essere divisi in gruppi e la prova svolta a turni. I candidati
portano un documento di identità e una penna, e devono conoscere il proprio
indirizzo di posta istituzionale e il numero di matricola.

La prova valuta entrambe le parti, studenti e docenti: serve a sapere dove sta
andando il corso. Chiude i primi tre blocchi (Tabella~\ref{tab:intro:midterm});
il quarto, la macchina, si esamina all'orale. Negli esempi della prova
compaiono istruzioni \riscv: il materiale del primo contatto del
Capitolo~\ref{ch:firstcontact} si dà per noto, e il codice di \vermone, parti
da 1 a 3, è una fonte naturale di domande. Ogni capitolo di questo libro
finisce con domande dello stesso tipo, e la preparazione migliore è rispondere
settimana per settimana.

\begin{table}[htb]
\centering\small
\caption{Che cosa copre la prova intercorso.}\label{tab:intro:midterm}
\begin{tabularx}{\textwidth}{@{}llL@{}}
\toprule
\textbf{Blocco} & \textbf{Lezioni} & \textbf{Argomenti}\\
\midrule
Informazione & da 3 a 6 & storia della computazione; informazione digitale; rappresentazione dei numeri: interi e reali\\
Logica & da 8 a 11 & algebra di Boole; elettronica e porte logiche; macchine combinatorie\\
Stato e memoria & da 13 a 16 & macchine sequenziali: latch, flip-flop, macchine a stati; registri e bus; memoria e memoria permanente\\
\emph{Non nella prova} & da 18 a 24 & il processore, istruzioni e interruzioni, la scheda video, ingresso e uscita, le prospettive: parte dell'esame orale\\
\bottomrule
\end{tabularx}
\end{table}

\begin{prgwarn}[Niente tecnologia. Niente GPT.]
La prova intercorso e l'esame sono \emph{\#nogpt}. Non è ammesso alcun aiuto
tecnologico o umano di alcun tipo: niente dispositivi, niente ChatGPT o altri
modelli linguistici, niente IA generativa. Copiare annulla subito la prova;
comportamenti più gravi hanno conseguenze più gravi.
\end{prgwarn}

\subsection{Valutazione}

L'esame ha due parti, ed entrambe vanno superate.

\begin{itemize}
\item \textbf{Teoria (50\%)}: un esame orale con domande estemporanee su tutto
  il programma; alcune risposte richiedono uno schema o un calcolo alla
  lavagna. L'orale può anche abbassare il voto ottenuto nella prova
  intercorso.
\item \textbf{Pratica (50\%)}: l'esame di laboratorio con la prof.ssa Staffa.
\end{itemize}

Il voto finale è la media del voto di teoria e del voto di pratica. I risultati
della prova intercorso e degli esami valgono entro l'anno accademico corrente.

Per verbalizzare l'esame serve un voto positivo sia in teoria sia in pratica, e
la verbalizzazione avviene solo durante un appello di teoria. L'esame di
laboratorio si tiene di solito una settimana prima dell'appello di teoria, ed
entrambi si prenotano su \url{https://uniparthenope.esse3.cineca.it}. Due
domande vengono poste ogni anno. \emph{Si può sostenere l'esame di teoria prima
di aver superato quello pratico?} Sì: il risultato vale per un anno accademico,
ma l'esame viene verbalizzato solo dopo aver superato la parte pratica.
\emph{Ci si può prenotare all'appello di teoria senza la prova pratica?} Sì, il
caso è lo stesso.

\subsection{La settimana della prova e l'evento finale}

Durante la settimana della prova intercorso le lezioni sono sostituite da
seminari sulle architetture di calcolo emergenti. Alla fine del corso, uno
studente o un gruppo che presenta il proprio gioco alla cerimonia di chiusura,
una variante di \vermone o un gioco originale, può ottenere punti bonus per
l'esame. La mattinata conclusiva ha due metà: l'ultima lezione, oltre un solo
core, da \riscv superscalare ai supercalcolatori (Capitolo~\ref{ch:beyond}), e
una visita ai sistemi di supercalcolo del Dipartimento di Scienze e Tecnologie.
I giochi finiti si possono pubblicare sul Cartridge Store di \prg, come mostra
il Capitolo~\ref{ch:store}.

% ----------------------------------------------------------------------------
\section{Materiale}

\begin{description}[style=nextline,leftmargin=1.5em]
\item[Slide.] Ogni lezione ha due presentazioni: la presentazione della
  lezione, di al più venti slide, per i concetti, e la presentazione del
  codice, con ogni listato, il programma completo della lezione e la guida per
  provarlo su \prgqt, su QEMU e sulla scheda. Sono sempre disponibili in una
  cartella condivisa di Google Drive e vengono pubblicate periodicamente su
  \url{https://elearning.uniparthenope.it}. Sono generate da sorgente:
  correzioni e nuovi esercizi sono benvenuti come contributi.
\item[Programmi.] Il programma di ogni lezione è nella cartella \fname{code}
  del repository del corso,
  \url{https://github.com/riscv-prg32/computer-architecture-with-prg32}, con
  le presentazioni, questo libro e i test che eseguono ogni listato.
\item[Piattaforma.] \prg, libero e open source, su
  \url{https://github.com/riscv-prg32/PRG32}, con il suo libro di testo,
  \emph{Playing with RISC-V: Assembly, C, and Retro Games on the PRG32
  Platform} \citep{prg32textbook}.
\item[Libri di testo.] \citet{patterson2020}, \emph{Computer Organization and
  Design, RISC-V Edition} (in italiano: \emph{Struttura e progetto dei
  calcolatori}), e \citet{tanenbaum2013}, \emph{Structured Computer
  Organization} (in italiano: \emph{Architettura dei calcolatori. Un approccio
  strutturale}), sono i due riferimenti del corso. Il programma elenca anche
  \citet{stallings2021} e \citet{bucci2017}.
\item[Standard.] La documentazione ufficiale di \riscv su
  \url{https://riscv.org} \citep{riscvisa}.
\end{description}

Il materiale didattico vive in una cartella condivisa con gli studenti, che
possono contribuirvi; per ottenere l'accesso lo studente invia al docente
l'indirizzo di un account collegato al servizio. Ogni altra comunicazione con
il docente passa per l'indirizzo istituzionale dello studente.

\begin{prgaside}[Cineforum]
Trenta film e serie sull'informatica accompagnano il corso, per ricordare che
la materia ha una cultura oltre che una tecnica: \emph{The Social
Network} (2010), \emph{Steve Jobs} (2015), \emph{Pirates of Silicon Valley}
(1999), \emph{The Imitation Game} (2014), \emph{Hackers} (1995), \emph{Ex
Machina} (2014), \emph{Her} (2013), \emph{Tron} (1982), \emph{Tron: Legacy}
(2010), \emph{The Matrix} (1999), \emph{WarGames} (1983), \emph{Sneakers}
(1992), \emph{Jobs} (2013), \emph{Antitrust} (2001), \emph{Blackhat} (2015),
\emph{The Fifth Estate} (2013), \emph{The Circle} (2017), \emph{Silicon
Valley} (2014--2019), \emph{Source Code} (2011), \emph{Chappie} (2015),
\emph{Mr.\ Robot} (2015--2019), \emph{Ghost in the Shell} (1995),
\emph{Minority Report} (2002), \emph{Transcendence} (2014), \emph{Ready Player
One} (2018), \emph{A.I.\ Artificial Intelligence} (2001), \emph{Revolt}
(2017), \emph{The Thirteenth Floor} (1999), \emph{Eagle Eye} (2008) e
\emph{Terminator~2: Judgment Day} (1991). \emph{Ready Player One} è anche il
riferimento culturale del libro di \prg: la dimestichezza con gli strati
fondamentali è un'alfabetizzazione che conserva il suo valore.
\end{prgaside}

% ----------------------------------------------------------------------------
\section{Primo contatto: chi siete?}

La prima lezione finisce con delle domande rivolte all'aula, e un lettore che
lavora da solo può rivolgerle a sé stesso. Non servono a selezionare, ma a
capire da dove cominciare.

\begin{itemize}
\item \textbf{Hardware.} Avete mai assemblato un calcolatore? Avete giocato
  con un Arduino o un ESP32?
\item \textbf{Software.} Quanti linguaggi di programmazione conoscete? Avete
  mai letto una riga di assembly?
\item \textbf{Cultura.} Che cosa sono il retrocomputing e il retrogaming?
  Avete mai visitato un museo dell'informatica?
\item \textbf{Concetti.} Sapete dire la differenza tra ROM, RAM e memoria di
  massa? Che cos'è il cloud computing? E l'Internet delle cose?
\item \textbf{RISC-V.} Avete sentito parlare di \riscv? Quale processore c'è
  nel vostro telefono, nel vostro portatile, nel vostro router?
\end{itemize}

In aula le domande sono seguite da un breve questionario anonimo, il
\emph{Cappello Parlante}, su ciò che attrae di più ciascuno studente nel futuro
dell'informatica e della scienza; i risultati vengono discussi e usati per
calibrare esempi e progetti.

\section{Prima del prossimo capitolo}

Il Capitolo~\ref{ch:setup} installa tutto, passo per passo, su Windows, macOS
e Linux. Si va più in fretta se i download lunghi sono già fatti.

\begin{exercise}[preparazione]
\begin{enumerate}[nosep]
\item Procuratevi un portatile su cui vi sia permesso installare programmi.
\item Scaricate l'installer di \prgqt per il vostro sistema dalla pagina
  \emph{Releases} di \url{https://github.com/riscv-prg32/PRG32-QT}.
\item Con una buona connessione, avviate il download lungo: ESP-IDF, seguendo
  il \fname{README} di \url{https://github.com/riscv-prg32/PRG32} per il vostro
  sistema.
\item Se qualcosa fallisce, non combattete da soli: conservate il messaggio di
  errore.
\end{enumerate}
\expected{Cominciate il Capitolo~\ref{ch:setup} con i download fatti. In due
ore avete \prgqt in funzione e una cartuccia compilata sul vostro portatile.}
\end{exercise}

% ----------------------------------------------------------------------------
% Generated by tools/textbook/export_listings.js from decks/plan.js and code/l00.S: do not edit.
\section{Il programma di questo capitolo}\label{sec:lab:l00}

I listati di questo capitolo sono parti di un solo programma,
\fname{code/l00.S} nel repository del corso, stampato per intero nel
Listato~\ref{lst:lab:l00}. I test automatici del materiale del corso lo
compilano e lo eseguono sull'esecutore senza interfaccia di \prgqt{}, con
ingressi programmati, e controllano il frame; la
Figura~\ref{fig:lab:l00} è il frame scritto dal test.

\begin{description}[style=nextline,leftmargin=1.5em]
\item[Sullo schermo.] La parola ARC12 in bianco su uno schermo blu, e sotto HELLO, RISC-V.
\item[Che cosa premere.] Niente: il programma va da solo.
\end{description}

\begin{figure}[htb]
\centering
\includegraphics[width=0.5\textwidth]{lab_l00.png}
\caption{\fname{l00.S} su \prgqt{}.}\label{fig:lab:l00}
\end{figure}

\asmfile[caption={Il programma completo del capitolo, \fname{code/l00.S}.},label=lst:lab:l00]{code/l00/program.S}

Per eseguirlo su \prgqt{} i comandi sono quelli del
Listato~\ref{lst:lab:l00:run}; il Paragrafo~\ref{sec:prg:testing} dà gli
stessi passi per QEMU e per la scheda.

\begin{shell}[caption={Compilare \fname{l00.S} ed eseguirlo su \prgqt{}.},label=lst:lab:l00:run]
# 1. start PRG32-QT on this computer (it listens on port 8080)

# 2. in the PRG32 folder, with ESP-IDF active: build the cartridge
python3 -m prg32 cartridge build ~/arc12/code/l00.S --portable \
    --entry-prefix l00 --name l00 --out l00.prg32

# 3. send it to PRG32-QT and start it
python3 -m prg32 esp32c6 upload-and-run l00.prg32 \
    --url http://localhost:8080 --slot cart0

# 4. look, press, compare with the previous slide
\end{shell}


